\section{Justificativas de alinhamento}
\label{sec:justificativa}

Cada canal prioritário é avaliado pelos três critérios da atividade. A proposta de valor do
\projeto{} combina quatro elementos: \emph{prevenção automatizada}, \emph{orientação
confiável}, \emph{pagamento possível} (créditos e preços sociais) e \emph{histórico digital do
pet}. O cliente-alvo principal é o tutor de baixa renda (classes C, D e E), que adia
tratamentos por restrição financeira e recorre à automedicação, conforme as entrevistas.

% ---------------------------------------------------------------------
\subsection{Canais de comunicação}

\justcanal{Boca a boca comunitário e ONGs}
{A proposta depende de confiança: o tutor precisa acreditar que os créditos e os preços
sociais são reais. Quando ONGs e protetores recomendam o aplicativo, a mensagem de acesso
financeiro e orientação confiável chega já validada por quem o tutor conhece.}
{ONGs e protetores independentes atendem justamente famílias de baixa renda e tutores
engajados com o bem-estar animal. Grupos de WhatsApp de bairro são o meio de comunicação já
usado por esse perfil.}
{Custo monetário próximo de zero. Exige tempo da equipe para construir parcerias, mas cada ONG
funciona como multiplicador. É o canal de menor custo de aquisição no estágio atual.}

\justcanal{Pontos de atendimento parceiros}
{O tutor descobre o aplicativo no momento em que a dor aparece (consulta, vacina, dúvida), e a
mensagem ``pague de forma possível'' responde diretamente à necessidade daquele instante.}
{Clínicas de bairro, pet shops e agropecuárias são frequentados pelo tutor de baixa renda; o
QR code reduz o atrito do cadastro.}
{Custo limitado a impressão de cartazes. Interessa ao parceiro, que ganha novos clientes e
ocupação de horários ociosos, então a adesão tende a ser espontânea.}

\justcanal{Eventos presenciais}
{Campanhas de vacinação e feiras de adoção reforçam o valor de prevenção e permitem
demonstrar o aplicativo na prática, com cadastro do pet e geração imediata do calendário
vacinal.}
{Reúnem tutores já interessados em cuidado animal, inclusive os que nunca levaram o pet ao
veterinário, e permitem ajuda presencial a quem tem menos familiaridade digital.}
{Pode ser feito como convidado de eventos já existentes (poder público, ONGs), sem organizar
evento próprio. O custo é o tempo de uma pequena equipe de voluntários.}

% ---------------------------------------------------------------------
\subsection{Canais de distribuição/entrega}

\justcanal{Google Play Store}
{É o canal em que o aplicativo (calendário, triagem, carteira de créditos e prontuário)
chega ao tutor. Sem ele nenhum elemento da proposta é entregue.}
{Android predomina entre usuários de baixa renda no Brasil, e a primeira fase do produto é
100\% mobile justamente para facilitar a adesão imediata desse público.}
{A publicação é simples, de baixo custo, e o Flutter permite manter uma única base de código.
Faixas de teste fechado permitem validar o produto antes do lançamento amplo.}

\justcanal{Rede de clínicas e veterinários parceiros}
{É onde o valor final é entregue: o tutor troca créditos e preços sociais por atendimento
real. Sem uma rede local ativa, a promessa de pagar de forma possível não se cumpre.}
{Os parceiros são regionais e próximos ao tutor, o que reduz custo de deslocamento e reforça o
caráter comunitário da proposta.}
{O modelo é viável porque o prestador ganha ocupação de janelas ociosas e novos clientes. O
esforço é concentrado em poucos parceiros na fase inicial, com um piloto local.}

\justcanal{Cadastro e instalação assistidos}
{A primeira experiência define a retenção: ao sair do cadastro com o calendário preventivo do
pet já gerado, o tutor percebe o valor imediatamente.}
{Tutores com pouca familiaridade digital ou pouco espaço de armazenamento no celular
precisam de ajuda para instalar e cadastrar; o assistente presencial elimina essa barreira.}
{Realizado por voluntários e pela própria equipe durante eventos e visitas a parceiros, sem
custo adicional de estrutura.}

% ---------------------------------------------------------------------
\subsection{Canais de relacionamento}

\justcanal{Notificações push automatizadas}
{Lembretes de vacinação, vermifugação e retorno transformam o cuidado em rotina
acompanhada, atacando diretamente o esquecimento do calendário vacinal relatado nas
entrevistas.}
{Notificações no celular chegam sem custo ao tutor e não dependem de conexão contínua ou de
leitura de conteúdo longo.}
{Custo praticamente nulo com o Firebase Cloud Messaging, e o calendário é gerado
automaticamente a partir do cadastro do pet.}

\justcanal{Comunidade de tutores e banco de horas}
{É o diferencial em relação a um marketplace comum: a troca de serviços entre vizinhos gera
os créditos usados no atendimento veterinário e cria pertencimento e reputação.}
{Atende ao tutor voluntário (com tempo e pouco dinheiro) e ao tutor de baixa renda que quer
cuidar do pet sem desestabilizar o orçamento.}
{Custo operacional baixo, pois a rede é mantida pelos próprios usuários. Exige regras claras de
conversão de créditos e moderação inicial.}

\justcanal{Suporte humano assistido}
{Casos sensíveis (emergência, disputa de crédito, dúvida sobre triagem) exigem resposta
humana para preservar a confiança, além de reforçar que a triagem é orientativa e não
substitui o veterinário.}
{WhatsApp é o canal de mensagens mais familiar ao público-alvo e dispensa aprendizado de
nova ferramenta.}
{Pode começar com atendimento pela própria equipe em horário limitado, com respostas
prontas para as dúvidas mais frequentes, e ser ampliado conforme o volume.}

% ---------------------------------------------------------------------
\subsection{Matriz de alinhamento entre canais e proposta de valor}

A Tabela~\ref{tab:matriz} resume qualitativamente quais elementos da proposta cada canal
prioritário reforça ($\checkmark$ contribuição direta, $\circ$ contribuição indireta).

\begin{table}[htbp]
\centering
\small
\caption{Matriz de alinhamento entre os canais prioritários e a proposta de valor.}
\label{tab:matriz}
\begin{tabular}{L{4.6cm} C{2.1cm} C{2.1cm} C{2.5cm} C{2.1cm}}
\toprule
\cabfundo \cab{Canal} & \cab{Prevenção} & \cab{Orientação} & \cab{Pagamento possível} & \cab{Histórico digital} \\
\midrule
Boca a boca e ONGs & $\circ$ & $\checkmark$ & $\checkmark$ & $\circ$ \\
Pontos parceiros (QR) & $\circ$ & $\circ$ & $\checkmark$ & $\checkmark$ \\
Eventos presenciais & $\checkmark$ & $\checkmark$ & $\circ$ & $\checkmark$ \\
\midrule
Google Play Store & $\checkmark$ & $\checkmark$ & $\checkmark$ & $\checkmark$ \\
Clínicas parceiras & $\checkmark$ & $\circ$ & $\checkmark$ & $\checkmark$ \\
Cadastro assistido & $\checkmark$ & $\circ$ & $\circ$ & $\checkmark$ \\
\midrule
Notificações push & $\checkmark$ & $\circ$ & $\circ$ & $\circ$ \\
Comunidade e banco de horas & $\circ$ & $\circ$ & $\checkmark$ & \\
Suporte humano & & $\checkmark$ & $\checkmark$ & \\
\bottomrule
\end{tabular}
\end{table}
